%==============================================================================
%  Alpamayo 1.5 Performance Characterization -- Internship Technical Report
%
%  Build:  latexmk -pdf alpamayo_characterization_report.tex
%     or:  pdflatex alpamayo_characterization_report.tex  (run twice for refs)
%     or:  tectonic alpamayo_characterization_report.tex
%
%  Figures live in ./figures/ and are copied from the workspace results tree.
%  Source methodology: notes/characterization.md (last updated 2026-08-18).
%==============================================================================
\documentclass[11pt,a4paper]{article}

\usepackage[utf8]{inputenc}
\usepackage[T1]{fontenc}
\usepackage{lmodern}
\usepackage[margin=1in]{geometry}
\usepackage{graphicx}
\usepackage{booktabs}
\usepackage{array}
\usepackage{longtable}
\usepackage{multirow}
\usepackage{amsmath}
\usepackage{amssymb}
\usepackage{siunitx}
\usepackage{xcolor}
\usepackage{caption}
\usepackage{subcaption}
\usepackage{listings}
\usepackage{pifont}
\usepackage{enumitem}
\usepackage{fancyhdr}
\usepackage[hidelinks,colorlinks=true,linkcolor=black!70!blue,urlcolor=blue,citecolor=blue]{hyperref}

%------------------------------------------------------------------ house style
\graphicspath{{figures/}}
\sisetup{detect-weight=true, round-mode=places, round-precision=2, group-separator={,}}
\definecolor{codebg}{gray}{0.96}
\definecolor{accent}{RGB}{0,80,140}

\lstset{
  basicstyle=\ttfamily\footnotesize,
  backgroundcolor=\color{codebg},
  frame=single,
  framerule=0pt,
  rulecolor=\color{codebg},
  xleftmargin=6pt, xrightmargin=6pt,
  aboveskip=8pt, belowskip=8pt,
  breaklines=true,
  columns=fullflexible,
  keepspaces=true,
  showstringspaces=false,
}
% Inline monospace that tolerates underscores/specials in EVERY engine
% (pdflatex/xelatex/lualatex/tectonic). \detokenize re-stringifies the argument
% so `_` prints literally; T1 fontenc gives the correct underscore glyph. This
% replaces a \lstinline-based macro that broke inside tabular cells under pdflatex.
\newcommand{\code}[1]{\texttt{\small\detokenize{#1}}}

\pagestyle{fancy}
\fancyhf{}
\lhead{\small Alpamayo 1.5 Characterization}
\rhead{\small\leftmark}
\cfoot{\thepage}
\renewcommand{\headrulewidth}{0.4pt}

\captionsetup{font=small,labelfont=bf}
\setlist{itemsep=2pt,topsep=3pt}

% Keep figures near their sections instead of clustering at the end.
\renewcommand{\topfraction}{0.9}
\renewcommand{\bottomfraction}{0.8}
\renewcommand{\textfraction}{0.07}
\renewcommand{\floatpagefraction}{0.7}
\setcounter{topnumber}{3}
\setcounter{totalnumber}{5}

% Common column type for numeric tables.
\newcolumntype{R}{>{\raggedleft\arraybackslash}p{1.9cm}}

%==============================================================================
\begin{document}

%------------------------------------------------------------------ title page
\begin{titlepage}
\centering
\vspace*{1.5cm}
{\Large\scshape Internship Technical Report\par}
\vspace{0.4cm}
{\color{accent}\rule{\linewidth}{1.2pt}}
\vspace{0.9cm}

{\huge\bfseries Performance Characterization of the\\[0.25em]
Alpamayo~1.5 End-to-End Driving Model\\[0.35em]
on NVIDIA Jetson~AGX~Thor\par}
\vspace{0.5cm}
{\large A latency, energy, memory, and accuracy study of a\\
vision--language trajectory planner deployed in Autoware\par}
\vspace{0.6cm}
{\color{accent}\rule{\linewidth}{1.2pt}}

\vspace{1.6cm}
{\large\bfseries Author:~\rule{5cm}{0.4pt}\par}      % <-- fill in name
\vspace{0.3cm}
{\normalsize University of Michigan \quad$\cdot$\quad \texttt{teadem@umich.edu}\par}

\vfill
{\large Host platform: NVIDIA Jetson AGX Thor (Blackwell, SM110)\par}
{\large Model: Alpamayo~1.5 (10B) \quad$\cdot$\quad Runtime: ROS~2 / Autoware\par}
\vspace{0.8cm}
{\large Report date: September~2026\par}
{\normalsize\itshape Methodology frozen 2026-08-18; experiments extended and status refreshed 2026-09-02\par}
\end{titlepage}

%------------------------------------------------------------------ abstract + toc
\pagenumbering{roman}
\begin{abstract}
\noindent
This report characterizes the on-device inference behavior of \emph{Alpamayo~1.5},
a 10-billion-parameter vision--language driving model that consumes multi-camera
imagery and ego odometry and emits a planned ego trajectory for the Autoware
autonomy stack. All measurements were collected by replaying recorded driving
logs through the production ROS~2 node on an NVIDIA Jetson~AGX~Thor and
instrumenting the node for per-inference timing, Jetson rail power, analytical
FLOP counts, KV-cache/GPU-memory footprint, and (for Waymo) ego trajectory
accuracy against replayed odometry. Two datasets are used throughout: a
four-camera \textbf{Waymo} moving segment (16 images per inference) and a
one-camera \textbf{MCity} segment (4 images per inference). Under default
runtime parameters the model runs at roughly \SI{3}{\second}--\SI{4}{\second}
per inference on Waymo and \SI{2.0}{\second}--\SI{2.6}{\second} on MCity,
draws a mean GPU-rail power of \SIrange{20}{33}{\watt}, and spends the large
majority of its GPU time in memory-bound autoregressive reason-token decode
(\SI{71}{\percent} of stage time). We present controlled sweeps over the four
principal runtime levers---\emph{diffusion steps}, \emph{trajectory samples
(K)}, \emph{numeric precision}, and \emph{reasoning-token cap}---and show that
trajectory-sample count and precision dominate the latency/energy budget, while
the token cap is essentially inert because generated reasoning length rarely
approaches the cap. FP16 is found to be both slower and higher-energy than the
BF16 default on this Thor stack. Ego ADE/FDE against replayed odometry stays
sub-metre across all runnable Waymo configurations. A controlled same-dataset
camera-count ablation shows camera count is the strongest accuracy lever
(1$\to$4 cameras cuts ADE $3.6\times$ at $2.9\times$ energy), and a best-of-K
analysis shows the model's unpublished samples are $40$--$50\%$ more accurate
than the trajectory it actually emits. The report documents each analysis with
its methodology, exact run configuration, and result plots, and closes with a
status audit and prioritized list of remaining work.
\end{abstract}

\newpage
\tableofcontents
\newpage
\pagenumbering{arabic}

%==============================================================================
\section{Introduction}
\label{sec:intro}

Alpamayo~1.5 is an end-to-end driving model: rather than the classical
detect--track--predict--plan pipeline, a single vision--language model (VLM)
ingests the recent camera history and ego state, produces a chain-of-thought
reasoning trace, and hands a latent to a diffusion \emph{expert} that denoises a
trajectory over a small number of steps. The model is wrapped as a ROS~2 node
and publishes an Autoware \code{Trajectory} message, chain-of-thought text, and
visualization markers. The inference path is strictly serial---one GPU, one
inference in flight at a time---so end-to-end latency is the sum of four stages:
vision encode, VLM prefill, autoregressive reason-token decode, and
trajectory-diffusion expert.

\paragraph{Goal of this work.}
The purpose of the internship characterization was to answer, quantitatively and
reproducibly, \emph{how the model behaves as a deployed real-time system on
edge hardware}: how long each inference takes, how much energy and GPU memory it
costs, how those costs scale with the model's tunable runtime knobs, and whether
the planned trajectory remains accurate as those knobs are turned down for speed.
Every claim in this report is backed by an artifact under the workspace
\code{results/} tree, and every number is defined by an explicit formula in
Section~\ref{sec:methodology}.

\paragraph{Two operating rules.}
The study follows two standing rules that make the numbers comparable:
\begin{enumerate}
  \item \textbf{Dataset rule.} Every non-ground-truth characterization is run for
  \emph{both} Waymo and MCity when the data and harness exist. Ground-truth
  dependent accuracy (ADE/FDE, actor collision/clearance) remains Waymo-only
  until a trusted MCity future-trajectory target is added.
  \item \textbf{Default-parameter rule.} A sweep varies \emph{only} the parameter
  under test; all other model/runtime parameters stay at the current node
  defaults: \code{precision_mode=bf16}, \code{max_generation_length=64},
  \code{num_trajectory_samples=1}, \code{num_diffusion_steps=5}. Dataset-specific
  I/O (camera topics, odometry topic, replay rate, QoS) is held at its
  dataset default.
\end{enumerate}

\subsection{Platform and model under test}
\label{sec:platform}

\begin{table}[h]
\centering
\caption{System under test.}
\label{tab:sut}
\small
\begin{tabular}{@{}ll@{}}
\toprule
\textbf{Component} & \textbf{Configuration} \\
\midrule
Accelerator      & NVIDIA Jetson AGX Thor (Blackwell class, SM110) \\
Model            & Alpamayo 1.5, $\approx$10B parameters, BF16 weights \\
Reasoning tokenizer & \code{nvidia/Cosmos-Reason2-8B} (for FLOP token counting) \\
Runtime          & ROS 2 node (\code{alpamayo_node_loki}), single-inference worker \\
Expert           & Native PyTorch diffusion denoiser (default) \\
Attention        & Flash / memory-efficient attention, autocast BF16 \\
Processed images & $560\times1008$, 4 temporal frames per camera \\
Power sensing    & Jetson INA rails via \code{/sys/class/hwmon}, 200\,Hz sampler \\
\bottomrule
\end{tabular}
\end{table}

\subsection{Datasets}
\label{sec:datasets}

\begin{table}[h]
\centering
\caption{The two replay datasets. The Waymo-vs-MCity comparison is
simultaneously a camera-count comparison; see Section~\ref{sec:camera}.}
\label{tab:datasets}
\small
\begin{tabular}{@{}lccccc@{}}
\toprule
\textbf{Dataset} & \textbf{Cameras} & \textbf{Images/inf.} & \textbf{Odometry} & \textbf{Replay rate} & \textbf{GT accuracy} \\
\midrule
Waymo (moving) & 4 & 16 & replayed ego odom & $0.1\times$ & ego ADE/FDE (sanity) \\
MCity          & 1 & 4  & \code{/gps/odom}   & $0.5\times$ & N/A (no trusted GT) \\
\bottomrule
\end{tabular}
\end{table}

The Waymo segment carries a genuine future ego-odometry timeline, which we use as
a \emph{sanity} target for ADE/FDE (Section~\ref{sec:accuracy})---this is not the
official Waymo planning benchmark. MCity provides one compressed camera and
\code{/gps/odom} but no trusted future-trajectory target, so all MCity accuracy
figures are reported as N/A by design.

%==============================================================================
\section{Common Methodology and Metric Definitions}
\label{sec:methodology}

\subsection{End-to-end replay harness}
\label{sec:harness}
Waymo experiments run through \code{scripts/run_experiment.sh}; MCity through
\code{scripts/run_mcity_experiment.sh}. Each run records a \code{config.env}
provenance file, starts the power sampler, waits for the node to report the model
is ready, starts camera republishing, the FLOP logger, predicted-trajectory /
odometry logging, and the camera validator, optionally starts the stamped
navigation-text publisher, then plays the bag with \code{/clock} at the
configured rate. After drain, helper process groups are torn down and offline
post-processing computes energy and (Waymo) ego ADE/FDE. The per-run artifacts
that every downstream number is derived from are listed in
Table~\ref{tab:artifacts}.

\begin{table}[h]
\centering
\caption{Primary per-run artifacts produced by the harness.}
\label{tab:artifacts}
\small
\begin{tabular}{@{}lp{9.2cm}@{}}
\toprule
\textbf{Artifact} & \textbf{Meaning} \\
\midrule
\code{config.env}              & Dataset, token cap, diffusion steps, replay rate, node suffix, nav settings, params file. \\
\code{inference_metrics.csv}   & One row per completed inference: ROS + wall-clock stamps, stage/model durations, token counts, nav text, publication metadata. \\
\code{power_samples.csv}       & Jetson INA rail power at 200\,Hz, wall-clock stamped. \\
\code{energy_per_inference.csv}& Per-inference rail energy and average power, integrated over each model window. \\
\code{scene_flops.csv}         & Analytical per-inference FLOP estimate by stage and effective TFLOP/s. \\
\code{waymo_eval.csv}          & Per-prediction ego ADE/FDE against replayed future odometry. \\
\code{camera_validation.txt}, \code{camera_receipts.csv} & Camera counts, drops/duplicates, trajectory count, per-inference input freshness. \\
\code{kv_memory_metrics.csv}   & Per-inference KV tensor footprint and CUDA allocator peaks. \\
\bottomrule
\end{tabular}
\end{table}

\subsection{Formal metric definitions}
\label{sec:metrics}

\paragraph{Latency.} All node timestamps use wall-clock \code{time.time()} so
they can be joined against power samples. The reported model duration for
inference $i$ is the model window
\begin{equation}
T_i = \texttt{end\_wall\_s}_i - \texttt{model\_start\_wall\_s}_i,
\end{equation}
while \code{duration_s} additionally includes CPU/GPU input preparation. Per-stage
latency is reported only where NVTX/stage instrumentation exists for that run.

\paragraph{GPU energy and power.} Using the GPU rail power
$P_\mathrm{gpu}(t)=\texttt{vdd\_gpu\_w}(t)$ over the model window $[a_i,b_i]$,
\begin{equation}
E_{\mathrm{gpu},i} = \int_{a_i}^{b_i} P_\mathrm{gpu}(t)\,dt,
\qquad
\bar{P}_{\mathrm{gpu},i} = \frac{E_{\mathrm{gpu},i}}{b_i-a_i},
\end{equation}
computed by trapezoidal integration with linear interpolation at the window
edges. Net energy subtracts an idle baseline $P_\mathrm{idle}$ estimated from
samples outside all guarded inference windows:
\begin{equation}
E^{\mathrm{net}}_{\mathrm{gpu},i} = \int_{a_i}^{b_i} \max\!\big(P_\mathrm{gpu}(t)-P_\mathrm{idle},\,0\big)\,dt .
\end{equation}
Peak GPU power for an inference is $\max_t P_\mathrm{gpu}(t)$ inside the model
window; because it is sampler-cadence dependent, the sampler frequency
(200\,Hz) is reported alongside every peak.

\paragraph{Trajectory accuracy (ADE/FDE).} For prediction $i$ with $K_i$ matched
future waypoints, predicted points $p_{i,k}$ and aligned target points $g_{i,k}$
in a common local frame,
\begin{equation}
\mathrm{ADE}_i = \frac{1}{K_i}\sum_{k=1}^{K_i}\lVert p_{i,k}-g_{i,k}\rVert_2,
\qquad
\mathrm{FDE}_i = \lVert p_{i,K_i}-g_{i,K_i}\rVert_2 .
\end{equation}
Reported ADE/FDE are the arithmetic means of $\mathrm{ADE}_i,\mathrm{FDE}_i$
over predictions with enough aligned future points.

\paragraph{Analytical FLOP / achieved throughput.} The analytical operation
model (from \code{scripts/flop_model.py}) counts transformer matmuls with
multiply-add as 2~FLOPs and excludes layernorm, softmax, RoPE, bias, and
residual adds:
\begin{equation}
\mathrm{decode\ TFLOP} = R\cdot 0.0171,
\qquad
\mathrm{total\ TFLOP} = 63.81 + \mathrm{decode\ TFLOP},
\end{equation}
where $R$ is the generated reasoning-token count and $63.81$~TFLOP is the fixed
vision+prefill+expert baseline calibrated to the Waymo 4-camera prompt
($P=3331$ prompt tokens, 17.1\,GFLOP per decode token). Achieved throughput is
obtained only after joining to measured latency:
$\mathrm{TFLOP/s}_i = \mathrm{total\ FLOP}_i / T_i$. The baseline constant is
\emph{not} camera-count-correct for MCity, so MCity FLOP is reported as an
analytical logger result, not a normalized operation count.

\paragraph{KV cache and CUDA memory.} KV bytes are summed directly from the
model's \code{past_key_values} tensors after VLM generation and before diffusion
consumes them:
\begin{equation}
\mathrm{KV\ MiB}_i = \frac{1}{2^{20}}\sum_{j\in\text{unique tensors}} \texttt{numel}_j\cdot\texttt{element\_size}_j .
\end{equation}
Whole-inference memory pressure is measured separately as the CUDA allocator peak
delta,
$\Delta_i = (\text{peak allocated during }i - \text{allocated at start}_i)/2^{20}$,
which includes KV growth \emph{plus} activations, temporaries, and diffusion/expert
working memory---it is not KV-only.

\paragraph{Throughput and freshness.}
\begin{equation}
\mathrm{inference\ FPS} = \frac{\text{completed inferences}}{\text{elapsed wall/sim time}},
\qquad
\mathrm{trajectory\ FPS} = \frac{\text{published trajectories}}{\text{elapsed wall/sim time}} .
\end{equation}
A lower inference count than camera-frame count is expected because the node has
a warmup/history requirement and an active-inference guard that prevents
overlapping model executions when model latency exceeds the trigger period.

%==============================================================================
\section{Canonical Baselines}
\label{sec:baseline}

\subsection*{Methodology}
The canonical baselines are archived clean replays that fix all runtime knobs and
serve as the reference operating point for every later sweep. The Waymo baseline
is the moving segment replayed with a \code{_loki} node suffix, 4 cameras and 16
images per inference; the MCity baseline is the one-camera harness run with
\code{/gps/odom}. Both were captured at \code{NUM_DIFF_STEPS=2}; later
default-parameter work uses the node default of 5 steps, so these are the
\emph{clean diff-2} baselines, not default-parameter baselines and not the full
Waymo planning benchmark.

\subsection*{Configuration}
\begin{table}[h]
\centering
\caption{Baseline run configurations.}
\label{tab:baseline-cfg}
\small
\begin{tabular}{@{}lcc@{}}
\toprule
\textbf{Parameter} & \textbf{Waymo baseline} & \textbf{MCity baseline} \\
\midrule
Source dir & \code{results/exp_eval/moving_diff2} & \code{results/kv/mcity_diff2} (+ legacy) \\
Cameras / images & 4 / 16 & 1 / 4 \\
\code{max_generation_length} & 64 & 64 \\
\code{num_diffusion_steps} & 2 & 2 \\
\code{num_trajectory_samples} & 1 & 1 \\
Replay rate & $0.1\times$ & $\approx 0.5\times$ (backlog-safe) \\
\bottomrule
\end{tabular}
\end{table}

\subsection*{Results}
Table~\ref{tab:baseline-res} summarizes the two baselines; the paired bar summary
and latency/energy comparison are shown in Figures~\ref{fig:baseline-table}
and~\ref{fig:latency-energy}. The Waymo baseline completes 46 inferences at
\SI{3452}{\milli\second} mean model duration, draws \SI{30.8}{\watt} mean GPU
power and \SI{106}{\joule} GPU energy per inference, and holds mean ego ADE/FDE
of \SI{0.107}{\metre}/\SI{0.194}{\metre}. MCity, with a quarter of the image
tokens, is faster (\SI{2640}{\milli\second}) and lower-power (\SI{20.2}{\watt},
\SI{52.5}{\joule}) but generates \emph{more} reasoning tokens on average
($13.2$ vs $7.8$), reflecting scene/route differences rather than a camera
effect.

\begin{table}[h]
\centering
\caption{Canonical baseline metrics (mean unless noted; $\pm$ is std).}
\label{tab:baseline-res}
\small
\begin{tabular}{@{}lcc@{}}
\toprule
\textbf{Metric} & \textbf{Waymo (diff2)} & \textbf{MCity (diff2)} \\
\midrule
Completed inferences        & 46 & 133 \\
Model duration (ms)         & $3452.5\pm316.3$ & $2640.0\pm925.1$ \\
Reasoning tokens (mean / max) & $7.85$ / $20$ & $13.24$ / $53$ \\
Mean GPU power (W)          & $30.81\pm1.13$ & $20.23\pm1.26$ \\
GPU energy / inf. (J)       & $106.07\pm5.94$ & $52.52\pm14.82$ \\
Total board energy / inf. (J) & $261.0\pm18.6$ & $169.1\pm54.3$ \\
Total analytical TFLOP      & $63.94$ & $\approx 64$ (Waymo-calibrated) \\
Ego ADE / FDE mean (m)      & $0.107$ / $0.194$ & N/A \\
Ego ADE / FDE median (m)    & $0.031$ / $0.074$ & N/A \\
\bottomrule
\end{tabular}
\end{table}

\begin{figure}[htbp]
\centering
\includegraphics[width=0.92\linewidth]{baseline_summary_table.png}
\caption{Cross-dataset baseline summary table (rendered from
\code{characterization_*.csv}).}
\label{fig:baseline-table}
\end{figure}

\begin{figure}[htbp]
\centering
\includegraphics[width=0.85\linewidth]{latency_energy_bars.png}
\caption{Waymo vs MCity mean model latency and per-inference GPU energy. MCity's
smaller prompt lowers both latency and energy despite longer reasoning traces.}
\label{fig:latency-energy}
\end{figure}

%==============================================================================
\section{Latency Analysis}
\label{sec:latency-analysis}

This section characterizes how end-to-end inference latency---and the GPU
energy that tracks it almost linearly---responds to the model's principal
runtime knobs. It first fixes the baseline latency, throughput, and input
freshness, then sweeps each latency-relevant parameter in turn: the
reasoning-token cap, the number of diffusion steps, and numeric precision,
holding all other parameters at their node defaults. The trajectory-sample
count $K$---a fourth latency lever---is treated in its own section
(Section~\ref{sec:ksweep}) because its primary payload is the best-of-$K$
accuracy analysis.

\subsection{Baseline latency, throughput, and input freshness}
\label{sec:latency}

\subsubsection*{Methodology}
Latency is taken from the wall-clock model window $T_i$ in
\code{inference_metrics.csv} and compared frame-to-frame, dataset-to-dataset, and
across matched sweep groups. Sustained throughput is completed inferences (and
published trajectories) over the replay span; because replay is slowed
($0.1\times$/$0.5\times$) both a \emph{wall} FPS and a \emph{sim} FPS are
reported. Input freshness joins \code{camera_receipts.csv} to
\code{inference_metrics.csv} by inference index and reports
$\texttt{odom\_minus\_camera\_stamp}=\texttt{ros\_stamp\_s}-\texttt{last\_selected\_stamp\_s}$
(age of the selected camera history) and
$\texttt{camera\_staleness}=\texttt{last\_received\_stamp\_s}-\texttt{last\_selected\_stamp\_s}$
(whether newer frames existed but were not selected).

\subsubsection*{Configuration}
\begin{table}[htbp]
\centering
\caption{Latency / throughput / freshness measurement configuration.}
\label{tab:latency-cfg}
\small
\begin{tabular}{@{}lcc@{}}
\toprule
\textbf{Parameter} & \textbf{Waymo} & \textbf{MCity} \\
\midrule
Cameras / images per inf. & 4 / 16 & 1 / 4 \\
\code{max_generation_length} & 64 & 64 \\
\code{num_diffusion_steps} & 2 & 2 \\
\code{num_trajectory_samples} & 1 & 1 \\
\code{precision_mode} & bf16 & bf16 \\
Replay rate & $0.1\times$ & $0.5\times$ \\
Trigger period & \SI{0.1}{\second} & \SI{0.1}{\second} \\
Power sampler & \SI{200}{\hertz} & \SI{200}{\hertz} \\
Swept parameter & \multicolumn{2}{c}{none (baseline operating point)} \\
\bottomrule
\end{tabular}
\end{table}

\subsubsection*{Results}
Table~\ref{tab:throughput} gives sustained throughput at the baseline operating
point. Because model latency ($\sim$\SI{3.5}{\second} Waymo) exceeds the
$0.1\,\mathrm{s}$ trigger period even after the $0.1\times$ slowdown, the node's
active-inference guard drops the completed-inference count well below the camera
frame count---this is expected serial-inference behavior, not frame loss.
Figure~\ref{fig:throughput-tokens} shows the paired throughput/token bars, and
Figure~\ref{fig:latency-complexity} relates per-inference latency to scene
complexity (reasoning-token count), the only strongly latency-correlated
per-frame quantity.

\begin{table}[h]
\centering
\caption{Sustained throughput at the baseline operating point.}
\label{tab:throughput}
\small
\begin{tabular}{@{}lccc@{}}
\toprule
\textbf{Run} & \textbf{Completed inf.} & \textbf{Sim inference FPS} & \textbf{Wall inference FPS} \\
\midrule
Waymo baseline (\code{exp_eval/moving_diff2}) & 46 & 2.49 & 0.25 \\
MCity baseline (\code{kv/mcity_diff2})        & 137 & 0.53 & 0.26 \\
\bottomrule
\end{tabular}
\end{table}

\begin{figure}[htbp]
\centering
\includegraphics[width=0.85\linewidth]{throughput_tokens_bars.png}
\caption{Paired throughput and reasoning-token bars for Waymo and MCity.}
\label{fig:throughput-tokens}
\end{figure}

\begin{figure}[htbp]
\centering
\includegraphics[width=0.8\linewidth]{latency_vs_complexity.png}
\caption{Per-inference model latency vs scene complexity (reasoning-token count).
Decode length is the dominant per-frame latency driver, consistent with the
stage profile in Section~\ref{sec:nsight}.}
\label{fig:latency-complexity}
\end{figure}

\paragraph{Input freshness.}
Table~\ref{tab:freshness} reports the two per-inference freshness quantities of
Section~\ref{sec:metrics}, aggregated over the baseline reruns. The selected
camera history is effectively synchronous with the trajectory/odometry stamp:
$\texttt{odom\_minus\_camera\_stamp}$ is exactly \SI{0}{\second} on Waymo (camera
and odometry share replay stamps) and within \SI{-64}{\milli\second} on MCity
(camera marginally newer, with one $-2.3$\,s outlier). Crucially,
$\texttt{camera\_staleness}=0$ everywhere---the node always selected the newest
available frames, so no usable camera data was left unconsumed. The selection
buffer never starved (minimum 4 selected frames, buffer depth up to 12).

\begin{table}[htbp]
\centering
\caption{Input freshness at the baseline operating point (per-inference means).}
\label{tab:freshness}
\small
\begin{tabular}{@{}lcc@{}}
\toprule
\textbf{Quantity} & \textbf{Waymo} & \textbf{MCity} \\
\midrule
Inferences aggregated                & 43 & 137 \\
\code{odom_minus_camera_stamp} mean (s) & 0.000 & $-0.064$ \\
\code{odom_minus_camera_stamp} min / max (s) & 0.000 / 0.000 & $-2.296$ / $0.019$ \\
\code{camera_staleness} mean / max (s) & 0.000 / 0.000 & 0.000 / 0.000 \\
Selected frames / inf.\ (min)        & 4 & 4 \\
Buffer depth (max)                   & 12 & 12 \\
\bottomrule
\end{tabular}
\end{table}

\subsection{Effect of the maximum generation length (reasoning-token cap)}
\label{sec:tokencap}

\subsubsection*{Methodology}
The token-cap sweep varies \code{MAX_GEN_LEN} while holding replay segment,
diffusion steps, camera setup, node code, and replay rate fixed, using the
clean-only Waymo artifacts. Both the configured cap and the \emph{actual}
generated reasoning length are reported; a flat trend is expected if actual
tokens do not approach the cap.

\subsubsection*{Configuration}
\begin{table}[htbp]
\centering
\caption{Reasoning-token-cap sweep configuration (clean-only Waymo).}
\label{tab:tokencap-cfg}
\small
\begin{tabular}{@{}ll@{}}
\toprule
\textbf{Parameter} & \textbf{Value} \\
\midrule
Swept: \code{max_generation_length} & $\{16,\,32,\,64\}$ \\
\code{num_diffusion_steps}   & 5 \\
\code{num_trajectory_samples}& 1 \\
\code{precision_mode}        & bf16 \\
Cameras / images per inf.    & 4 / 16 (Waymo) \\
Replay rate                  & $0.1\times$ \\
Accuracy target              & ego-odometry sanity ADE/FDE only \\
\bottomrule
\end{tabular}
\end{table}

\subsubsection*{Results}
The cap is effectively inert: mean generated reasoning length is
$8.1/9.5/9.5$ tokens against caps of $16/32/64$, never hitting the cap
(max-token hit rate $0\%$ across all runs). Consequently latency
($3958/3963/4007$\,ms) and GPU energy ($112.6/113.2/113.6$\,J) are flat within
noise (Table~\ref{tab:tokencap}). ADE/FDE drift slightly \emph{downward} with a
larger cap ($0.031\to0.028$\,m ADE), but the change is inside run-to-run
variance. This confirms that on these segments the reasoning-token cap is not a
useful latency lever---diffusion steps and K are. Figures~\ref{fig:tokencap-le},
\ref{fig:tokencap-acc}, and~\ref{fig:tokencap-tflops} show the sweep.

\begin{table}[h]
\centering
\caption{Clean-only Waymo token-cap sweep (mean).}
\label{tab:tokencap}
\small
\begin{tabular}{@{}lcccccc@{}}
\toprule
\textbf{Cap} & \textbf{$n$} & \textbf{Actual tokens} & \textbf{Latency (ms)} & \textbf{GPU E (J)} & \textbf{ADE (m)} & \textbf{FDE (m)} \\
\midrule
16 & 12 & 8.1 & $3957.5$ & $112.6$ & 0.0314 & 0.0726 \\
32 & 12 & 9.5 & $3962.5$ & $113.2$ & 0.0302 & 0.0664 \\
64 & 12 & 9.5 & $4006.5$ & $113.6$ & 0.0275 & 0.0627 \\
\bottomrule
\end{tabular}
\end{table}

\begin{figure}[htbp]
\centering
\begin{subfigure}{0.49\linewidth}
  \centering
  \includegraphics[width=\linewidth]{token_cap_latency_energy.png}
  \caption{Latency / power / energy vs cap.}
  \label{fig:tokencap-le}
\end{subfigure}\hfill
\begin{subfigure}{0.49\linewidth}
  \centering
  \includegraphics[width=\linewidth]{token_cap_accuracy.png}
  \caption{Ego ADE/FDE vs cap.}
  \label{fig:tokencap-acc}
\end{subfigure}
\caption{Token-cap sweep: cost and accuracy are both flat because generated
length never reaches the cap.}
\label{fig:tokencap}
\end{figure}

\begin{figure}[htbp]
\centering
\includegraphics[width=0.7\linewidth]{token_cap_tflops.png}
\caption{Analytical TFLOP and achieved TFLOP/s across token caps
($\approx$\SI{16}{TFLOP/s}, memory-bound).}
\label{fig:tokencap-tflops}
\end{figure}

\subsubsection*{Dataset-matched MCity token-cap sweep}
\label{sec:tokencap-mcity}
The clean token-cap sweep above was Waymo-only; the dataset rule asks for MCity
too. We ran the matched MCity sweep (\code{MAX_GEN_LEN}$\in\{16,32,64\}$,
\code{NUM_DIFF_STEPS=5}, \code{K=1}, bf16, rate $0.5\times$;
\code{results/token_cap_mcity/}). The result differs qualitatively from Waymo:
mean reasoning length \emph{rises} with the cap ($7.5\!\to\!10.3\!\to\!13.4$
tokens, Table~\ref{tab:tokencap-mcity}), because MCity's natural reasoning
distribution extends past 16/32 tokens, so the lower caps genuinely truncate it
(the maximum observed length sits right at the cap: 15, 25, 46). The cap is
therefore a \emph{mild active constraint} on MCity, unlike the inert Waymo case.
It still does not produce a clean latency/energy trend---decode is a small share
of MCity latency, and each cap completes a different subset of inferences over the
replay (138/136/176), so scene mix dominates the modest differences. Note the
analytical TFLOP column reports the live-logger value ($\approx$64), which for
MCity is the Waymo-calibrated over-count; the camera-normalized MCity baseline is
\SI{22.65}{TFLOP} (Section~\ref{sec:flops-mcity}). Figure~\ref{fig:tokencap-mcity}
shows the sweep.

\begin{table}[htbp]
\centering
\caption{Dataset-matched MCity token-cap sweep (mean). Accuracy is N/A (no MCity GT).}
\label{tab:tokencap-mcity}
\small
\begin{tabular}{@{}lccccc@{}}
\toprule
\textbf{Cap} & \textbf{$n$} & \textbf{Tokens (mean / max)} & \textbf{Latency (ms)} & \textbf{GPU W} & \textbf{GPU E (J)} \\
\midrule
16 & 138 & 7.54 / 15 & 2286.8 & 18.55 & 42.30 \\
32 & 136 & 10.33 / 25 & 2485.5 & 19.09 & 46.77 \\
64 & 176 & 13.40 / 46 & 2148.2 & 22.00 & 46.07 \\
\bottomrule
\end{tabular}
\end{table}

\begin{figure}[htbp]
\centering
\includegraphics[width=\linewidth]{token_cap_mcity_sweep.png}
\caption{MCity token-cap sweep: reasoning tokens, latency, GPU power, and GPU
energy. Reasoning length responds to the cap (unlike Waymo), but the cost metrics
are dominated by per-cap scene mix.}
\label{fig:tokencap-mcity}
\end{figure}

\subsection{Effect of the number of diffusion steps}
\label{sec:diffusion}

\subsubsection*{Methodology}
The diffusion-step sweep varies \code{NUM_DIFF_STEPS} while holding token cap,
trajectory samples, precision, replay segment, camera setup, node code, and
replay rate fixed \emph{within} each dataset. Because diffusion changes
expert-loop work directly, it is expected to separate latency/energy more cleanly
than the token cap when reasoning length stays low.

\subsubsection*{Configuration}
\begin{table}[htbp]
\centering
\caption{Diffusion-step sweep configuration.}
\label{tab:diffusion-cfg}
\small
\begin{tabular}{@{}lcc@{}}
\toprule
\textbf{Parameter} & \textbf{Waymo} & \textbf{MCity} \\
\midrule
Swept: \code{num_diffusion_steps} & \multicolumn{2}{c}{$\{2,\,5,\,10\}$} \\
\code{max_generation_length}  & 64 & 64 \\
\code{num_trajectory_samples} & 1 & 1 \\
\code{precision_mode}         & bf16 & bf16 \\
Cameras / images per inf.     & 4 / 16 & 1 / 4 \\
Replay rate                   & $0.1\times$ & $0.5\times$ \\
\bottomrule
\end{tabular}
\end{table}

\subsubsection*{Results}
Table~\ref{tab:diffusion} shows monotonic MCity growth in latency and energy with
steps (\SI{2005}{\milli\second}$\to$\SI{2643}{\milli\second};
\SI{47}{\joule}$\to$\SI{57}{\joule}). On Waymo the trend is present but partly
masked by decode-length variance, since the fixed 4-camera prefill/decode
dominates the step-scaled expert work. Waymo ego ADE/FDE is essentially flat
across steps (ADE $0.104$--$0.117$\,m), so more diffusion steps buy little
trajectory accuracy at this operating point. Figure~\ref{fig:diffusion} shows the
full sweep.

\begin{table}[h]
\centering
\caption{Diffusion-step sweep (mean). MCity accuracy is N/A by design.}
\label{tab:diffusion}
\small
\begin{tabular}{@{}llccccc@{}}
\toprule
\textbf{Dataset} & \textbf{Steps} & \textbf{$n$} & \textbf{Latency (ms)} & \textbf{GPU E (J)} & \textbf{ADE (m)} & \textbf{FDE (m)} \\
\midrule
\multirow{3}{*}{Waymo}
 & 2  & 45 & $3551.9$ & $107.6$ & 0.109 & 0.205 \\
 & 5  & 46 & $3498.3$ & $112.1$ & 0.117 & 0.188 \\
 & 10 & 43 & $3663.0$ & $118.6$ & 0.104 & 0.201 \\
\midrule
\multirow{3}{*}{MCity}
 & 2  & 200 & $2005.1$ & $47.0$ & N/A & N/A \\
 & 5  & 156 & $2260.1$ & $50.2$ & N/A & N/A \\
 & 10 & 131 & $2642.7$ & $56.9$ & N/A & N/A \\
\bottomrule
\end{tabular}
\end{table}

\begin{figure}[htbp]
\centering
\includegraphics[width=0.95\linewidth]{diffusion_step_sweep.png}
\caption{Dataset-matched diffusion-step sweep: latency, GPU energy/power,
reasoning tokens, and (Waymo) ADE/FDE vs \code{NUM_DIFF_STEPS}.}
\label{fig:diffusion}
\end{figure}

\subsection{Effect of quantization (numeric precision)}
\label{sec:precision}

\subsubsection*{Methodology}
The precision study runs a full replay for each viable precision mode, varying
\emph{only} \code{precision_mode} and comparing latency, GPU power/energy, KV/CUDA
memory, and (Waymo) accuracy. FP8/INT8 are not run end-to-end because no
end-to-end Alpamayo FP8/INT8 runtime exists in this repo; the preserved torchao
INT8/FP8 VLM attempt is a negative micro-benchmark (it falls back to a slow
dequant-to-BF16 path) and is explicitly excluded from replay metrics.

\subsubsection*{Configuration}
\begin{table}[htbp]
\centering
\caption{Precision-replay sweep configuration.}
\label{tab:precision-cfg}
\small
\begin{tabular}{@{}lcc@{}}
\toprule
\textbf{Parameter} & \textbf{Waymo} & \textbf{MCity} \\
\midrule
Swept: \code{precision_mode} & \multicolumn{2}{c}{bf16, fp16 \; (fp8/int8 unsupported)} \\
\code{num_diffusion_steps}   & 5 & 5 \\
\code{max_generation_length} & 64 & 64 \\
\code{num_trajectory_samples}& 1 & 1 \\
\code{QUANT}                 & none & none \\
Cameras / images per inf.    & 4 / 16 & 1 / 4 \\
Replay rate                  & $0.1\times$ & $0.5\times$ \\
\bottomrule
\end{tabular}
\end{table}

\subsubsection*{Results}
The headline result is counter-intuitive but robust: \textbf{FP16 is both slower
and higher-energy than the BF16 default on this Thor stack}, and roughly doubles
KV and CUDA-peak memory (Table~\ref{tab:precision}). Waymo FP16 runs
\SI{7531}{\milli\second} vs BF16's \SI{4010}{\milli\second} and burns
\SI{221.8}{\joule} vs \SI{116.6}{\joule} GPU per inference; MCity shows the same
$\approx$2.7$\times$ latency and $3\times$ energy penalty. Trajectory accuracy is
statistically unchanged (Waymo ADE $0.128$ vs $0.130$\,m), so FP16 offers no
accuracy upside to justify its cost. FP8/INT8 remain unsupported for full replay.
Figure~\ref{fig:precision} shows the comparison bars.

\begin{table}[h]
\centering
\caption{Default-parameter precision replay (mean). FP16 is uniformly worse than
BF16 here.}
\label{tab:precision}
\small
\begin{tabular}{@{}llcccccc@{}}
\toprule
\textbf{Dataset} & \textbf{Mode} & \textbf{$n$} & \textbf{Lat.\ (ms)} & \textbf{GPU E (J)} & \textbf{KV (MiB)} & \textbf{CUDA$\Delta$ (MiB)} & \textbf{ADE (m)} \\
\midrule
\multirow{2}{*}{Waymo} & BF16 & 41 & $4009.5$ & $116.6$ & 433.7 & 972.5 & 0.128 \\
                       & FP16 & 23 & $7530.9$ & $221.8$ & 868.3 & 1523.7 & 0.130 \\
\midrule
\multirow{2}{*}{MCity} & BF16 & 163 & $2334.0$ & $52.0$ & 119.1 & 268.2 & N/A \\
                       & FP16 & 73  & $6211.7$ & $161.7$ & 237.8 & 478.2 & N/A \\
\midrule
\multicolumn{2}{@{}l}{FP8 / INT8} & 0 & \multicolumn{5}{c}{unsupported end-to-end runtime (see \code{*_fp8/UNSUPPORTED.md})} \\
\bottomrule
\end{tabular}
\end{table}

\begin{figure}[htbp]
\centering
\includegraphics[width=0.92\linewidth]{precision_bars.png}
\caption{BF16 vs FP16 latency, GPU power, GPU energy, and (Waymo) accuracy.}
\label{fig:precision}
\end{figure}

%==============================================================================
\section{Latency Breakdown by Stage}
\label{sec:nsight}

\subsection*{Methodology}
End-to-end latency is the serial sum of four GPU stages---vision encode, VLM
prefill, autoregressive reason decode, and the trajectory-diffusion expert. To
attribute latency to each, the node installs NVTX stage ranges
(\code{profile_stages:=true}) and an \code{nsys} capture is sliced per stage with
\code{nvtx_gpu_proj_sum}. Kernel-name aggregation alone cannot separate stages
that share GEMM kernels, so the measured split relies on the NVTX ranges.
Captures are the single-camera MCity/staged configuration; per-stage
\emph{energy} is not available unless power samples are aligned to the same stage
timestamps (listed as remaining work).

\subsection*{Results}
The stage profile (Table~\ref{tab:stages}) explains every latency observation in
this report: \textbf{autoregressive reason decode is \SI{71.4}{\percent} of GPU
time}, which is why decode length is the dominant per-frame latency driver
(Section~\ref{sec:latency-analysis}). Prefill is \SI{12.8}{\percent}, the
trajectory-diffusion expert \SI{10.2}{\percent} (the step-scaled work that the
diffusion and $K$ sweeps move), and vision encode only \SI{5.6}{\percent}.
Figure~\ref{fig:perframe} shows the per-frame stage time series.

\begin{table}[h]
\centering
\caption{GPU-time stage breakdown (NVTX, single-camera staged capture).}
\label{tab:stages}
\small
\begin{tabular}{@{}lccp{6.2cm}@{}}
\toprule
\textbf{Stage} & \textbf{GPU time (s)} & \textbf{Share} & \textbf{Character} \\
\midrule
\code{reason_decode}         & 55.16 & 71.4\% & Memory-bound per-token decoder GEMMs; the latency driver. \\
\code{prefill}               & 9.91  & 12.8\% & Compute-dense full-sequence GEMMs + flash attention. \\
\code{traj_diffusion_expert} & 7.86  & 10.2\% & Expert 2-CTA cluster GEMMs; scales with steps and K. \\
\code{vision_encode}         & 4.34  & 5.6\%  & ViT flash attention + biased GEMMs + heavy pointwise. \\
\bottomrule
\end{tabular}
\end{table}

\begin{figure}[htbp]
\centering
\includegraphics[width=0.85\linewidth]{per_frame_timeseries.png}
\caption{Per-frame stage timing series from the staged NVTX capture.}
\label{fig:perframe}
\end{figure}

%==============================================================================
\section{Bottleneck Identification with Nsight Profiling}
\label{sec:bottleneck}

\subsection*{Methodology}
Beyond the stage shares, an Nsight Systems capture with CUDA + NVTX + OSRT traces
and CUDA memory-usage tracking is used to inspect the dominant kernels inside each
stage and classify whether each stage is compute- or memory-bound
(Table~\ref{tab:nsight-cfg}). This is what turns the ``where does time go'' split
of Section~\ref{sec:nsight} into a ``what is the bottleneck'' answer that guides
optimization.

\subsection*{Configuration}
\begin{table}[htbp]
\centering
\caption{Nsight profiling configuration.}
\label{tab:nsight-cfg}
\small
\begin{tabular}{@{}ll@{}}
\toprule
\textbf{Parameter} & \textbf{Value} \\
\midrule
Profiler        & Nsight Systems (\code{nsys}) \\
Trace domains   & \code{cuda}, \code{nvtx}, \code{osrt} \\
Stage ranges    & NVTX, \code{profile_stages:=true} \\
Aggregation     & \code{nvtx_gpu_proj_sum} (measured stage split) \\
Capture         & single-camera MCity / staged \\
Memory tracking & \code{--cuda-memory-usage=true} \\
\bottomrule
\end{tabular}
\end{table}

\subsection*{Results}
The bottleneck is unambiguous: the \SI{71}{\percent}-of-time reason-decode stage
is \textbf{memory-bandwidth-bound}, dominated by small-tile $M{=}1$ \code{nvjet}
GEMMs that stream the decoder's weights once per generated token
(Table~\ref{tab:kernels}). This is exactly why the model sustains only
$\approx$\SI{16}{TFLOP/s} despite $\approx$\SI{64}{TFLOP} of analytical work
(Section~\ref{sec:flops})---it is starved on weight bandwidth, not compute. By
contrast, prefill is \textbf{compute-dense} (a few large full-sequence GEMMs plus
full flash attention), and the diffusion expert is a thin block of 2-CTA cluster
GEMMs plus memory-efficient attention---a fusion candidate. The optimization
implication is direct: latency is won by \emph{shrinking the memory-bound decode}
(fewer reasoning tokens, quantized weights, smaller KV) rather than by adding
FLOPs, and by fusing the expert loop---not by touching the already-cheap vision
and prefill stages.

\begin{table}[htbp]
\centering
\caption{Dominant kernels and boundness per stage (staged \code{nsys} capture).}
\label{tab:kernels}
\small
\begin{tabular}{@{}lp{6.4cm}l@{}}
\toprule
\textbf{Stage} & \textbf{Dominant kernels} & \textbf{Boundness} \\
\midrule
\code{reason_decode} & small-tile $M{=}1$ \code{nvjet} decoder QKV/MLP GEMMs (per token), KV concat & memory-bandwidth \\
\code{prefill} & few large full-sequence \code{nvjet} GEMMs + full flash attention & compute \\
\code{traj_diffusion_expert} & 2-CTA cluster \code{nvjet} GEMMs + \code{fmha_cutlass} mem-efficient attention & mixed (fusion candidate) \\
\code{vision_encode} & ViT flash attention + biased GEMMs + heavy pointwise & mixed \\
\bottomrule
\end{tabular}
\end{table}

%==============================================================================
\section{Power and Energy Analyses}
\label{sec:power}

\subsection*{Methodology}
\code{power_sampler.py} reads Jetson AGX Thor INA sensors under
\code{/sys/class/hwmon} at 200\,Hz, resolving rails by sensor name. The key rails
are \code{total_w} (VDD\_IN board input), \code{vdd_gpu_w} (GPU rail, used for
GPU-only reporting), \code{vdd_cpu_soc_mss_w}, and \code{vin_sys_5v0_w}.
\code{energy_per_inference.py} joins power to inferences by wall-clock time and
integrates each rail over $[a_i,b_i]$; GPU joules per inference use
\code{vdd_gpu_w_energy_j}, and net energy subtracts the idle baseline over the
active duration. Energy is reported per inference, not per stage (stage-level
energy needs NVTX-aligned power integration and is listed as remaining work).

\subsection*{Configuration}
\begin{table}[htbp]
\centering
\caption{GPU power / energy integration configuration.}
\label{tab:power-cfg}
\small
\begin{tabular}{@{}ll@{}}
\toprule
\textbf{Parameter} & \textbf{Value} \\
\midrule
Reported GPU rail        & \code{vdd_gpu_w} \\
Integration window       & model window $[\texttt{model\_start\_wall\_s},\texttt{end\_wall\_s}]$ \\
Integration method       & trapezoidal, linear interpolation at window edges \\
Idle baseline            & median rail power outside guarded windows \\
Net energy               & $\int \max(P_\mathrm{gpu}-P_\mathrm{idle},0)\,dt$ \\
Power sampler rate       & \SI{200}{\hertz} \\
Datasets                 & Waymo ($0.1\times$) and MCity ($0.5\times$), baseline + sweeps \\
\bottomrule
\end{tabular}
\end{table}

\subsection*{Results}
GPU-rail energy tracks latency almost linearly across configurations because mean
GPU power is nearly constant within a dataset (\SIrange{20}{33}{\watt}). At the
baseline, Waymo draws \SI{106}{\joule} GPU / inference vs MCity's
\SI{52.5}{\joule}. Figure~\ref{fig:mcity-rail} decomposes MCity energy by rail,
and Figure~\ref{fig:mcity-timepower} shows the GPU power time series with the
integrated per-inference windows.

\begin{figure}[htbp]
\centering
\begin{subfigure}{0.49\linewidth}
  \centering
  \includegraphics[width=\linewidth]{mcity_energy_by_rail.png}
  \caption{MCity per-inference energy by rail.}
  \label{fig:mcity-rail}
\end{subfigure}\hfill
\begin{subfigure}{0.49\linewidth}
  \centering
  \includegraphics[width=\linewidth]{mcity_gpu_energy_time_power.png}
  \caption{MCity GPU power time series and integrated windows.}
  \label{fig:mcity-timepower}
\end{subfigure}
\caption{MCity power/energy decomposition. GPU is a minority of total board
power; CPU/SOC and 5\,V system rails contribute substantially to board energy.}
\label{fig:power}
\end{figure}

\subsection*{Peak GPU power}
For each inference we also take the maximum \code{vdd_gpu_w} sample inside the
model window and summarize it across inferences (mean / p95 / peak). Peaks are
sampler-cadence dependent, so the \SI{200}{\hertz} sampler frequency is reported
and peaks from logs with differing cadence are never compared
(Table~\ref{tab:peak-cfg}).

\begin{table}[htbp]
\centering
\caption{Peak GPU power configuration.}
\label{tab:peak-cfg}
\small
\begin{tabular}{@{}ll@{}}
\toprule
\textbf{Parameter} & \textbf{Value} \\
\midrule
Reported GPU rail   & \code{vdd_gpu_w} \\
Per-inference statistic & $\max_t P_\mathrm{gpu}(t)$ inside the model window \\
Aggregation         & mean / p95 / peak across inferences \\
Runs                & diffusion-matched (Waymo + MCity, steps $\{2,5,10\}$) \\
Power sampler rate  & \SI{200}{\hertz} \\
\bottomrule
\end{tabular}
\end{table}

Active peak GPU power is bounded well below the module's total board peak: on
Waymo it reaches $\approx$\SI{52}{\watt} (e.g. \code{waymo_diff10} active peak
\SI{52.3}{\watt}) and on MCity $\approx$\SI{40}{\watt} (\code{mcity_diff2}
\SI{40.1}{\watt}), against total-board peaks of \SIrange{123}{125}{\watt}.
The GPU rail is therefore \emph{not} the board-power bottleneck; CPU/SOC and
system rails dominate total draw. Figure~\ref{fig:peak} shows the paired peak-GPU
bars.

\begin{figure}[htbp]
\centering
\includegraphics[width=0.8\linewidth]{peak_power_bars.png}
\caption{Peak GPU-rail power, Waymo vs MCity (200\,Hz sampler).}
\label{fig:peak}
\end{figure}

%==============================================================================
\section{Predicted Trajectory Accuracy}
\label{sec:accuracy}

\subsection*{Methodology}
Two independent accuracy views are used, both Waymo-only.
\textbf{(a) Ego ADE/FDE sanity.} \code{evaluate_waymo_predictions.py} groups
predicted trajectory points by prediction, builds a de-duplicated odometry
timeline, sets each prediction's evaluation start to
$\texttt{prediction\_stamp}+\texttt{stamp\_offset\_s}$, interpolates odometry at
each $\texttt{time\_from\_start}$, transforms future poses into the
evaluation-start local frame, and computes the 2D displacement errors of
Section~\ref{sec:metrics}. This is an ego-odometry sanity check, not the official
Waymo planning benchmark.
\textbf{(b) GT actor interaction.} \code{evaluate_waymo_gt_actors.py} loads
Waymo GT actor boxes, represents the planned ego as a $4.8\times2.1$\,m oriented
rectangle, and uses separating-axis polygon intersection for collisions and
polygon-edge distance for clearance---reporting collisions, near-misses
($<1.0$\,m), and minimum clearance. It does \emph{not} compute ADE/FDE.

\subsection*{Configuration}
\begin{table}[htbp]
\centering
\caption{Trajectory-accuracy evaluation configuration (Waymo baseline).}
\label{tab:accuracy-cfg}
\small
\begin{tabular}{@{}ll@{}}
\toprule
\textbf{Parameter} & \textbf{Value} \\
\midrule
Dataset / operating point & Waymo baseline, \code{num_diffusion_steps=2} \\
Ego-odometry target       & replayed future odometry (sanity, not GT benchmark) \\
Planned-ego footprint     & oriented rectangle $4.8 \times 2.1$\,m \\
GT actor classes          & vehicle, pedestrian, cyclist \\
Actor \code{max_range_m}  & 120 \\
\code{gt_tolerance_s}     & 0.06 \\
\code{near_miss_m}        & 1.0 \\
\code{actor_prefilter_m}  & 80.0 \\
\bottomrule
\end{tabular}
\end{table}

\subsection*{Results}
Over 46 baseline predictions the model holds mean ego ADE/FDE of
\SI{0.107}{\metre}/\SI{0.194}{\metre} (median \SI{0.031}{\metre}/\SI{0.074}{\metre})
against replayed odometry---sub-metre and stable, as expected for a short-horizon
sanity target. The GT actor check reports \textbf{zero collisions} across the
segment, with the nearest annotated actor a vehicle at $\sim$\SI{0.68}{\metre}
minimum clearance. Figure~\ref{fig:accuracy} shows the per-prediction error and
summary bars.

\begin{figure}[htbp]
\centering
\begin{subfigure}{0.56\linewidth}
  \centering
  \includegraphics[width=\linewidth]{accuracy_per_prediction.png}
  \caption{Per-prediction ADE/FDE.}
  \label{fig:accuracy-pp}
\end{subfigure}\hfill
\begin{subfigure}{0.42\linewidth}
  \centering
  \includegraphics[width=\linewidth]{accuracy_bars.png}
  \caption{Aggregate ADE/FDE.}
  \label{fig:accuracy-bars}
\end{subfigure}
\caption{Waymo ego trajectory accuracy against replayed odometry (baseline
diff-2 operating point).}
\label{fig:accuracy}
\end{figure}

%==============================================================================
\section{FLOP / TOPS}
\label{sec:flops}

\subsection*{Methodology}
The FLOP logger subscribes to \code{/alpamayo/reasoning_stamped}, tokenizes the
reasoning text with \code{nvidia/Cosmos-Reason2-8B}, and applies the analytical
model of Section~\ref{sec:metrics}: a fixed \SI{63.81}{TFLOP} vision+prefill+expert
baseline (calibrated to the Waymo 4-camera, $P=3331$-token prompt) plus
\SI{17.1}{GFLOP} per decode token. Achieved TFLOP/s is
$\mathrm{total\ FLOP}/T_i$. Absolute values are analytical estimates suitable for
\emph{relative/stage} analysis, not hardware-counter truth; MCity's baseline term
is not camera-normalized and is reported as a logger result only.

\subsection*{Configuration}
\begin{table}[htbp]
\centering
\caption{Analytical FLOP model configuration.}
\label{tab:flop-cfg}
\small
\begin{tabular}{@{}ll@{}}
\toprule
\textbf{Parameter} & \textbf{Value} \\
\midrule
Fixed baseline FLOP   & \SI{63.81}{TFLOP} (Waymo 4-camera prompt, $P=3331$ tokens) \\
Per-decode-token cost & \SI{17.1}{GFLOP} \\
Reasoning tokenizer   & \code{nvidia/Cosmos-Reason2-8B} \\
Matmul convention     & multiply-add counted as 2 FLOPs \\
Excluded operations   & layernorm, softmax, RoPE, bias, residual add \\
Achieved throughput   & $\mathrm{total\ FLOP}/T_i$ (joined to measured latency) \\
\bottomrule
\end{tabular}
\end{table}

\subsection*{Results}
Because generated reasoning length is small ($\sim$8--13 tokens), the decode term
(\SI{0.15}{TFLOP}) is a rounding error on the fixed \SI{63.81}{TFLOP} baseline, so
total analytical FLOP is effectively constant at $\approx$\SI{64}{TFLOP} per
inference regardless of token cap. Achieved throughput on the clean Waymo
token-cap runs is only $\approx$\SI{16.3}{TFLOP/s}
(Table~\ref{tab:flops})---the model is far from compute-bound, consistent with the
memory-bound decode profile of Section~\ref{sec:nsight}.
Figures~\ref{fig:flops-tokens} and~\ref{fig:tokens-decode} show the per-frame FLOP
vs token relationship and the decode-token contribution.

\begin{table}[h]
\centering
\caption{Analytical FLOP and achieved throughput, clean Waymo token-cap runs.}
\label{tab:flops}
\small
\begin{tabular}{@{}lcccc@{}}
\toprule
\textbf{Run} & \textbf{Reason tokens} & \textbf{Decode TFLOP} & \textbf{Total TFLOP} & \textbf{Achieved TFLOP/s} \\
\midrule
\code{tok16} & 8.8 & 0.151 & 63.96 & 16.27 \\
\code{tok32} & 9.5 & 0.162 & 63.97 & 16.43 \\
\code{tok64} & 9.5 & 0.162 & 63.97 & 16.27 \\
\bottomrule
\end{tabular}
\end{table}

\begin{figure}[htbp]
\centering
\begin{subfigure}{0.49\linewidth}
  \centering
  \includegraphics[width=\linewidth]{flops_tokens_per_frame.png}
  \caption{Per-frame analytical FLOP vs reasoning tokens.}
  \label{fig:flops-tokens}
\end{subfigure}\hfill
\begin{subfigure}{0.49\linewidth}
  \centering
  \includegraphics[width=\linewidth]{tokens_vs_decode.png}
  \caption{Decode-token contribution to total FLOP.}
  \label{fig:tokens-decode}
\end{subfigure}
\caption{The fixed vision/prefill/expert baseline dominates total FLOP; decode
tokens add a negligible slope.}
\label{fig:flops}
\end{figure}

\begin{figure}[htbp]
\centering
\includegraphics[width=0.8\linewidth]{scene_flops_grouped.png}
\caption{Per-inference analytical FLOP grouped across replay reruns, illustrating
run-to-run stability of the operation-count estimate.}
\label{fig:scene-flops}
\end{figure}

\subsection*{Camera-normalized MCity FLOP}
\label{sec:flops-mcity}
The live logger tags \emph{every} run with the Waymo 4-camera baseline
(\SI{63.81}{TFLOP} fixed), which is not camera-correct for MCity's single camera.
Recomputing the fixed FLOP with only the camera-dependent inputs changed
($N_\mathrm{img}: 16\!\to\!4$, prompt $P: 3331\!\to\!1171$, i.e. 720 image
tokens $+$ 451 text) yields a much smaller MCity baseline
(Table~\ref{tab:flops-mcity}). The recomputation is validated by the Waymo column:
its 5-step fixed reproduces the logger's \SI{63.81}{TFLOP} constant exactly. At
the default 5 steps the true MCity fixed FLOP is \SI{22.65}{TFLOP}, so the
Waymo-calibrated logger \textbf{over-counts MCity operations by
$\mathbf{2.82\times}$}; the gap is dominated by prefill
(\SI{49.55}{TFLOP}$\to$\SI{16.67}{TFLOP}) as the prompt shrinks. MCity per-inference
FLOP/TOPS should therefore be reported against this normalized baseline, not the
Waymo constant. (Artifact: \code{scripts/flop_model_camera.py},
\code{results/flops_camera_normalized/summary.csv}.)

\begin{table}[htbp]
\centering
\caption{Camera-normalized analytical fixed FLOP per inference. Only the
camera-dependent inputs differ between rows; the Waymo 5-step value reproduces the
validated \SI{63.81}{TFLOP} logger constant.}
\label{tab:flops-mcity}
\small
\begin{tabular}{@{}lcccccc@{}}
\toprule
\textbf{Dataset} & \textbf{$N_\mathrm{img}$} & \textbf{$P$} & \textbf{Steps} & \textbf{Vision} & \textbf{Prefill} & \textbf{Fixed TFLOP} \\
\midrule
\multirow{3}{*}{Waymo (4 cam)}
 & 16 & 3331 & 2  & 10.50 & 49.55 & 61.55 \\
 & 16 & 3331 & 5  & 10.50 & 49.55 & \textbf{63.81} \\
 & 16 & 3331 & 10 & 10.50 & 49.55 & 67.58 \\
\midrule
\multirow{3}{*}{MCity (1 cam)}
 & 4 & 1171 & 2  & 2.63 & 16.67 & 20.64 \\
 & 4 & 1171 & 5  & 2.63 & 16.67 & \textbf{22.65} \\
 & 4 & 1171 & 10 & 2.63 & 16.67 & 26.01 \\
\bottomrule
\end{tabular}
\end{table}

%==============================================================================
\section{KV Cache Size}
\label{sec:kv}

\subsection*{Methodology}
KV bytes are summed from the model's \code{past_key_values} after VLM generation
and before diffusion consumes the cache; a recursive walker handles HuggingFace
cache objects, dicts/lists/tuples, and \code{key_cache}/\code{value_cache}
attributes while de-duplicating Python object IDs. CUDA peak allocation delta is
measured independently by resetting peak stats at inference start and reading peak
allocated/reserved after trajectory publication; it captures whole-inference
memory pressure (KV + activations + temporaries + diffusion working memory), not
KV alone.

\subsection*{Configuration}
\begin{table}[htbp]
\centering
\caption{KV / GPU-memory measurement configuration.}
\label{tab:kv-cfg}
\small
\begin{tabular}{@{}lcc@{}}
\toprule
\textbf{Parameter} & \textbf{Waymo} & \textbf{MCity} \\
\midrule
Cameras / images per inf.    & 4 / 16 & 1 / 4 \\
\code{max_generation_length} & 64 & 64 \\
\code{num_diffusion_steps}   & 2 & 2 \\
\code{num_trajectory_samples}& 1 & 1 \\
\code{precision_mode}        & bf16 & bf16 \\
CFG / navigation             & off (unguided KV $=0$) & off \\
KV measurement point         & \multicolumn{2}{c}{after VLM generation, before diffusion} \\
\bottomrule
\end{tabular}
\end{table}

\subsection*{Results}
KV footprint scales with prompt length, which scales with camera count: Waymo's
$\sim$3084-token prompt yields \SI{433.6}{MiB} mean KV vs MCity's 845-token prompt
at \SI{118.9}{MiB} (Table~\ref{tab:kv}). Yet KV is a small fraction of
whole-inference memory pressure---$1.96\%$ of the Waymo CUDA peak and $0.55\%$ of
MCity's---so the CUDA peak allocation delta (\SI{972.5}{MiB} Waymo,
\SI{268.2}{MiB} MCity) is dominated by activations and diffusion/expert working
memory, not the cache. The trajectory-start token was decoded on $100\%$ of Waymo
and $94\%$ of MCity inferences. Figures~\ref{fig:kv-bars} and~\ref{fig:kv-scatter}
show the paired KV bars and the KV-vs-tokens scatter.

\begin{table}[h]
\centering
\caption{KV cache and CUDA memory (baseline reruns, mean).}
\label{tab:kv}
\small
\begin{tabular}{@{}lcc@{}}
\toprule
\textbf{Metric} & \textbf{Waymo} & \textbf{MCity} \\
\midrule
Prompt KV sequence length     & 3083.7 & 845.2 \\
Total KV (MiB)                & 433.6 (max 437.1) & 118.9 (max 125.3) \\
Model parameters (GiB)        & 20.64 & 20.64 \\
CUDA peak allocated (GiB)     & 21.63 & 20.92 \\
CUDA peak alloc.\ delta (MiB) & 972.5 & 268.2 \\
KV \% of peak allocated       & 1.96\% & 0.55\% \\
Trajectory-start hit rate     & 1.00 & 0.94 \\
\bottomrule
\end{tabular}
\end{table}

\begin{figure}[htbp]
\centering
\begin{subfigure}{0.49\linewidth}
  \centering
  \includegraphics[width=\linewidth]{kv_memory_bars.png}
  \caption{Waymo vs MCity KV and CUDA peak.}
  \label{fig:kv-bars}
\end{subfigure}\hfill
\begin{subfigure}{0.49\linewidth}
  \centering
  \includegraphics[width=\linewidth]{kv_memory_scatter.png}
  \caption{KV footprint vs token count.}
  \label{fig:kv-scatter}
\end{subfigure}
\caption{KV cache footprint and its (small) share of total GPU memory pressure.}
\label{fig:kv}
\end{figure}

%==============================================================================
\section{Trajectory-Sample (K) Sweep}
\label{sec:ksweep}

\subsection*{Methodology}
This sweep varies \code{num_trajectory_samples} (K) while holding all other knobs
at defaults within each dataset. The Autoware \code{Trajectory} output publishes
sample~0 only, so the reported Waymo ADE/FDE is \emph{sample-0} accuracy; a true
best-of-K $\mathrm{minADE}_K/\mathrm{minFDE}_K$ requires logging all returned
samples and is tracked as remaining work.

\subsection*{Configuration}
\begin{table}[htbp]
\centering
\caption{Trajectory-sample (K) sweep configuration.}
\label{tab:ksweep-cfg}
\small
\begin{tabular}{@{}lcc@{}}
\toprule
\textbf{Parameter} & \textbf{Waymo} & \textbf{MCity} \\
\midrule
Swept: \code{num_trajectory_samples} (K) & \multicolumn{2}{c}{$\{1,\,2,\,4\}$} \\
\code{max_generation_length} & 64 & 64 \\
\code{num_diffusion_steps}   & 5 & 5 \\
\code{precision_mode}        & bf16 & bf16 \\
Cameras / images per inf.    & 4 / 16 & 1 / 4 \\
Replay rate                  & $0.1\times$ & $0.5\times$ \\
Published sample             & \multicolumn{2}{c}{sample 0 (Autoware \code{Trajectory})} \\
\bottomrule
\end{tabular}
\end{table}

\subsection*{Results}
K is the single most expensive latency/energy lever measured: Waymo latency rises
$3010\to4662\to8474$\,ms and GPU energy $107\to190\to376$\,J as K goes
$1\to2\to4$ (Table~\ref{tab:ksweep}), i.e. close to linear in K because each
sample re-runs the diffusion expert. Sample-0 ADE actually \emph{worsens} with K
($0.103\to0.181$\,m), which is expected: adding samples does not improve the
published sample~0, and any best-of-K benefit is invisible until all samples are
scored. Figure~\ref{fig:ksweep} shows the full sweep.

\begin{table}[h]
\centering
\caption{Trajectory-sample sweep (mean). Waymo accuracy is sample-0 only.}
\label{tab:ksweep}
\small
\begin{tabular}{@{}llccccc@{}}
\toprule
\textbf{Dataset} & \textbf{K} & \textbf{$n$} & \textbf{Latency (ms)} & \textbf{GPU E (J)} & \textbf{ADE$_0$ (m)} & \textbf{FDE$_0$ (m)} \\
\midrule
\multirow{3}{*}{Waymo}
 & 1 & 53 & $3009.6$ & $106.6$ & 0.103 & 0.184 \\
 & 2 & 35 & $4661.9$ & $190.4$ & 0.138 & 0.181 \\
 & 4 & 21 & $8473.8$ & $375.7$ & 0.181 & 0.259 \\
\midrule
\multirow{3}{*}{MCity}
 & 1 & 185 & $2159.5$ & $51.5$ & N/A & N/A \\
 & 2 & 55  & $3261.9$ & $85.6$ & N/A & N/A \\
 & 4 & 89  & $4903.1$ & $148.2$ & N/A & N/A \\
\bottomrule
\end{tabular}
\end{table}

\begin{figure}[htbp]
\centering
\includegraphics[width=0.95\linewidth]{trajectory_sample_sweep.png}
\caption{Default-parameter K$=1/2/4$ sweep: latency, GPU power/energy, tokens, and
Waymo sample-0 ADE/FDE.}
\label{fig:ksweep}
\end{figure}

\subsection*{Best-of-K accuracy}
\label{sec:bestofk}
The sample-0 accuracy above answers ``is the \emph{published} path better with
more samples?'' (no). The more useful question is ``does the model's sample set
\emph{contain} a better path?''---i.e.\ $\mathrm{minADE}_K/\mathrm{minFDE}_K$. That
was previously unanswerable because only sample~0 is published. We added an
opt-in all-samples logger to the node (\code{all_samples_csv}), re-ran Waymo
K$=1/2/4$ logging every returned sample, and scored each sample against replayed
odometry with the same alignment as Section~\ref{sec:accuracy}
(\code{scripts/evaluate_best_of_k.py}). Table~\ref{tab:bestofk} and
Figure~\ref{fig:bestofk} show the result: \textbf{the published sample~0 is far
from the best sample.} Taking the best of $K$ cuts ADE by $41\%$ at K$=2$
(\SI{0.145}{\metre}$\to$\SI{0.085}{\metre}) and $48\%$ at K$=4$
(\SI{0.199}{\metre}$\to$\SI{0.104}{\metre}), with similar $37$--$49\%$ FDE
reductions. So the extra samples are not wasted work---their diversity holds
materially more accurate trajectories---but capturing that gain requires either
publishing/selecting among samples or a learned ranker, since sample~0 alone even
degrades slightly as K grows. (The per-K prediction counts differ---53/35/21 over
the segment---and K$=4$'s $n{=}21$ is small, so read the within-run
sample-0-vs-best gap, which is robust, rather than cross-K minADE.)

\begin{table}[htbp]
\centering
\caption{Best-of-K vs published sample~0 (Waymo, replayed-odometry accuracy).}
\label{tab:bestofk}
\small
\begin{tabular}{@{}lccccc@{}}
\toprule
\textbf{K} & \textbf{$n$} & \textbf{ADE$_0$ (m)} & \textbf{minADE$_K$ (m)} & \textbf{FDE$_0$ (m)} & \textbf{minFDE$_K$ (m)} \\
\midrule
1 & 53 & 0.115 & 0.115 & 0.234 & 0.234 \\
2 & 35 & 0.145 & \textbf{0.085} & 0.191 & \textbf{0.120} \\
4 & 21 & 0.199 & \textbf{0.104} & 0.310 & \textbf{0.159} \\
\bottomrule
\end{tabular}
\end{table}

\begin{figure}[htbp]
\centering
\includegraphics[width=0.85\linewidth]{best_of_k_bars.png}
\caption{Best-of-K accuracy: the published sample~0 (grey) vs the best of $K$
samples (color). The model's sample diversity contains substantially more
accurate trajectories than the one it publishes.}
\label{fig:bestofk}
\end{figure}

%==============================================================================
\section{Camera-Count Ablation}
\label{sec:cameraablation}

\subsection*{Methodology}
The Waymo-vs-MCity comparison of Section~\ref{sec:camera} confounds camera count
with dataset, scene, topic, and route. To isolate the camera-count effect we ran
a \emph{controlled} ablation: the same Waymo segment replayed three times with the
node consuming 1, 2, or 4 cameras---achieved by handing the node a subset of
\code{camera_topics}/\code{camera_indices} (front; front $+$ front-right; all
four) while every other parameter and the replayed data are held fixed. The
republisher still publishes all four compressed streams; only the node's
subscription set changes. The model accepts a variable camera set without error,
so both the cost curve and the accuracy curve are measurable.

\subsection*{Configuration}
\begin{table}[htbp]
\centering
\caption{Camera-count ablation configuration (same Waymo segment).}
\label{tab:cameraablation-cfg}
\small
\begin{tabular}{@{}ll@{}}
\toprule
\textbf{Parameter} & \textbf{Value} \\
\midrule
Swept: camera count & 1 (front) / 2 (front $+$ front-right) / 4 (all) \\
\code{camera_indices} & [1] / [1,\,2] / [0,\,1,\,2,\,6] \\
\code{max_generation_length} & 64 \\
\code{num_diffusion_steps}   & 5 \\
\code{num_trajectory_samples}& 1 \\
\code{precision_mode}        & bf16 \\
Replay rate                  & $0.1\times$ \\
\bottomrule
\end{tabular}
\end{table}

\subsection*{Results}
This is the cleanest cost/accuracy trade in the study
(Table~\ref{tab:cameraablation}, Figure~\ref{fig:cameraablation}): every cost
metric scales monotonically with camera count, and so does accuracy---in the
opposite direction. Going from 1 to 4 cameras \textbf{cuts ego ADE
$3.6\times$} (\SI{0.333}{\metre}$\to$\SI{0.092}{\metre}) and \textbf{FDE
$6\times$} (\SI{0.905}{\metre}$\to$\SI{0.152}{\metre}), but costs $2.0\times$
latency, $2.9\times$ GPU energy, and $3.7\times$ KV cache (the prompt grows
$839\to3083$ tokens). The 4-camera point reproduces the canonical baseline
(KV \SI{433.5}{MiB} vs \SI{433.6}{MiB}; ADE \SI{0.092}{\metre} vs the baseline
range), which validates the ablation harness. Two points worth flagging:
\begin{itemize}
  \item The single-camera model still plans a usable trajectory (sub-metre ADE)
  but its FDE nearly reaches \SI{1}{\metre}---the missing side/telephoto views
  cost long-horizon accuracy, as expected.
  \item The live FLOP logger reports $\approx$\SI{64}{TFLOP} for all three rows
  because it uses the fixed Waymo 4-camera constant; the true camera-normalized
  operation count scales with camera count (Section~\ref{sec:flops-mcity}:
  \SI{20.6}{TFLOP} at 1 camera vs \SI{63.8}{TFLOP} at 4), so the real compute
  saving of dropping cameras is larger than the logged column suggests.
\end{itemize}

\begin{table}[htbp]
\centering
\caption{Same-dataset camera-count ablation (mean). Only the camera count varies.}
\label{tab:cameraablation}
\small
\begin{tabular}{@{}lccccccc@{}}
\toprule
\textbf{Cameras} & \textbf{$n$} & \textbf{Lat.\ (ms)} & \textbf{GPU E (J)} & \textbf{KV (MiB)} & \textbf{Prompt tok} & \textbf{ADE (m)} & \textbf{FDE (m)} \\
\midrule
1 & 88 & 1468.9 & 34.5 & 117.9 & 839 & 0.333 & 0.905 \\
2 & 79 & 1852.0 & 54.6 & 223.0 & 1586 & 0.151 & 0.429 \\
4 & 53 & 2951.9 & 101.0 & 433.5 & 3083 & 0.092 & 0.152 \\
\bottomrule
\end{tabular}
\end{table}

\begin{figure}[htbp]
\centering
\includegraphics[width=\linewidth]{camera_ablation_sweep.png}
\caption{Camera-count ablation on a single Waymo segment: compute cost and KV rise
with camera count while ego ADE/FDE fall. This isolates the camera-count effect
that the cross-dataset comparison could not.}
\label{fig:cameraablation}
\end{figure}

%==============================================================================
\section{Token-Limit Behavior}
\label{sec:tokenlimit}

\subsection*{Methodology}
Token-limit behavior records, per run, the actual generated reasoning-token
count, whether the \code{max_generation_length} cap was reached, whether any
inference produced zero reasoning tokens, whether the trajectory-start token was
decoded, and whether the trajectory was truncated before a valid decode. Cap hits
are inferred as $\texttt{reason\_tokens}\ge\texttt{max\_generation\_length}$; the
trajectory-start hit rate is taken from the KV instrumentation reruns.

\subsection*{Configuration}
\begin{table}[htbp]
\centering
\caption{Token-limit-behavior data sources.}
\label{tab:tokenlimit-cfg}
\small
\begin{tabular}{@{}ll@{}}
\toprule
\textbf{Parameter} & \textbf{Value} \\
\midrule
Runs analyzed        & Waymo baseline (\code{max_generation_length=64}) + MCity baseline \\
Cap-hit rule         & $\texttt{reason\_tokens}\ge\texttt{max\_generation\_length}$ \\
Traj-start hit rate  & from KV reruns (generated-token slots) \\
Truncation markers   & token IDs not persisted (inferred, not explicit) \\
\bottomrule
\end{tabular}
\end{table}

\subsection*{Results}
At the default cap of 64 the model never comes close to truncation
(Table~\ref{tab:tokenlimit}): the maximum observed reasoning length is 20 tokens
on Waymo and 47 on MCity, so the max-token hit rate is $0\%$ on both datasets and
every inference decoded a complete 64-point trajectory (zero incomplete-trajectory
rows). The trajectory-start token was decoded on $100\%$ of Waymo and $94\%$ of
MCity inferences. MCity occasionally emits a zero-reasoning-token inference
(8 of 137, $\sim6\%$)---the model still produced a valid trajectory in those cases
---whereas Waymo never did. This confirms the finding of
Section~\ref{sec:tokencap}: the reasoning-token cap is not an active constraint at
the default operating point.

\begin{table}[htbp]
\centering
\caption{Token-limit behavior at the baseline operating point (cap $=64$).}
\label{tab:tokenlimit}
\small
\begin{tabular}{@{}lcc@{}}
\toprule
\textbf{Quantity} & \textbf{Waymo} & \textbf{MCity} \\
\midrule
Inferences                    & 46 & 137 \\
Reasoning tokens (mean / max) & 7.85 / 20 & 13.33 / 47 \\
Max-token hits (rate)         & 0 (0\%) & 0 (0\%) \\
Zero-reasoning-token rows     & 0 & 8 \\
Incomplete-trajectory rows    & 0 & 0 \\
Trajectory-start hit rate     & 1.00 & 0.94 \\
\bottomrule
\end{tabular}
\end{table}

%==============================================================================
\section{Camera Stream Validation and Freshness}
\label{sec:camera}

\subsection*{Methodology}
Camera delivery is measured two ways because \code{inference_metrics.csv} records
only \emph{configured} camera counts, not real receipts.
\code{validate_camera_streams.py} subscribes to the compressed topics the node
consumes (plus optional raw topics and \code{/alpamayo/predicted_trajectory}) and,
on shutdown, writes per-topic message counts, measured Hz, dropped-gap counts
(inter-frame gap $>1.5\times$ median dt), duplicate-stamp counts, cross-camera
skew, and trajectory count. Separately, \code{camera_receipts.csv} yields the two
freshness quantities of Section~\ref{sec:latency}. The Waymo-vs-MCity comparison
is confounded by camera count, dataset, topic, rate, and route, so it is
\emph{not} a controlled camera-count ablation (that is listed as remaining work).

\subsection*{Configuration}
\begin{table}[htbp]
\centering
\caption{Camera-validation and freshness configuration.}
\label{tab:camera-cfg}
\small
\begin{tabular}{@{}ll@{}}
\toprule
\textbf{Parameter} & \textbf{Value} \\
\midrule
Subscribed topics   & compressed camera topics $+$ \code{/alpamayo/predicted_trajectory} \\
Drop threshold      & inter-frame gap $> 1.5\times$ median dt \\
Reported metrics    & msg count, Hz, drops, duplicate stamps, cross-camera skew \\
Freshness metrics   & \code{odom_minus_camera_stamp}, \code{camera_staleness} \\
QoS probe matrix    & reliable vs.\ best-effort, node vs.\ no-node \\
\bottomrule
\end{tabular}
\end{table}

\subsection*{Results}
Under the clean baseline every camera delivers its full frame count at the
expected \SI{10}{\hertz} with zero drops, zero duplicates, and zero cross-camera
skew, and the node publishes a trajectory for every completed inference
(Table~\ref{tab:camval}). Heavier clean token-cap sweeps exposed intermittent
BEST\_EFFORT front-camera loss (e.g. \code{199/187/172/199}), which
reliable-QoS / no-node probes isolated to transport backpressure rather than
model latency. This motivates the standing recommendation to use reliable QoS for
characterization captures where transport fidelity matters.

\begin{table}[htbp]
\centering
\caption{Camera-stream validation at the baseline operating point.}
\label{tab:camval}
\small
\begin{tabular}{@{}lcc@{}}
\toprule
\textbf{Quantity} & \textbf{Waymo} & \textbf{MCity} \\
\midrule
Cameras                    & 4 & 1 \\
Messages / camera          & 199 (all four equal) & 2649 \\
Measured rate (Hz)         & 10.0 & 10.0 \\
Dropped-gap count          & 0 & 0 \\
Duplicate-stamp count      & 0 & 0 \\
Cross-camera stamp skew    & 0\,ms & n/a (single camera) \\
Predicted trajectories     & 43 & 137 \\
\bottomrule
\end{tabular}
\end{table}

%==============================================================================
\section{Navigation-Text Input}
\label{sec:nav}

\subsection*{Methodology}
Final navigation-text experiments use the \emph{stamped} input path, not static
parameter overrides. The harness starts \code{run_nav_text_publisher.sh} when
\code{NAV_TEXT}/\code{NAV_TEXT_SCHEDULE} is set; it publishes
\code{autoware_internal_debug_msgs/StringStamped} on \code{/alpamayo/nav_input}
aligned to \code{/clock}. The node selects the latest nav message whose stamp is
not newer than the odometry/input stamp for that inference, and records the
effective \code{nav_text} on each completed inference row, so a valid nav run must
show the intended text on completed rows.

\subsection*{Configuration}
\begin{table}[htbp]
\centering
\caption{Navigation-text input configuration.}
\label{tab:nav-cfg}
\small
\begin{tabular}{@{}ll@{}}
\toprule
\textbf{Parameter} & \textbf{Value} \\
\midrule
Input topic     & \code{/alpamayo/nav_input} \\
Message type    & \code{autoware_internal_debug_msgs/StringStamped} \\
Publish rate    & \SI{10}{\hertz} (\code{NAV_TEXT_HZ}) \\
Stamp alignment & aligned to \code{/clock} \\
Selection rule  & latest nav stamp not newer than the inference input stamp \\
Proof text      & ``stop, red light'' (stamped) \\
Method          & stamped input (static parameter override deprecated) \\
\bottomrule
\end{tabular}
\end{table}

\subsection*{Results}
We ran a controlled same-settings A/B/C: three Waymo replays identical in every
parameter (\code{MAX_GEN_LEN=64}, \code{NUM_DIFF_STEPS=5}, \code{K=1}, bf16, rate
$0.1\times$) differing only in the stamped nav text---\emph{A: ``drive straight
through the intersection''}, \emph{B: ``turn right at the intersection''},
\emph{C: ``stop, red light ahead''}. Each run recorded exactly its intended nav
text on every inference (confirming the stamped path end-to-end), and all three
went through the CFG-nav dual-pass model call. Two findings emerge
(Table~\ref{tab:nav}, Figure~\ref{fig:nav}):

\begin{enumerate}
  \item \textbf{Nav text barely moves the trajectory on this segment.} The
  pairwise trajectory divergence between prompts is tiny: mean endpoint
  divergence is \SI{1.4}{\centi\metre} (straight vs right), \SI{3.3}{\centi\metre}
  (straight vs stop), and \SI{18.7}{\centi\metre} (right vs stop), against
  $\sim$\SI{40}{\metre}-scale paths. Ego ADE stays sub-metre and similar across
  prompts (\SIrange{0.082}{0.106}{\metre}). On a moving segment where the road
  geometry strongly constrains the path, contradictory nav commands are largely
  overridden by the scene---an important caveat for anyone expecting nav text to
  steer the planner.
  \item \textbf{CFG-nav is not free.} The guided+unguided dual pass makes these
  runs $\sim$\SIrange{4.8}{5.9}{\second} per inference versus $\sim$\SI{3.0}{\second}
  for the non-nav \code{K=1} baseline (Section~\ref{sec:ksweep})---a $\sim$1.6$\times$
  latency and $\sim$1.7$\times$ GPU-energy cost for conditioning that changes the
  output only marginally here.
\end{enumerate}

\begin{table}[htbp]
\centering
\caption{Stamped-nav A/B/C (Waymo, same settings; only nav text differs).
Divergence is the mean trajectory / endpoint distance between prompt pairs at
matching prediction stamps.}
\label{tab:nav}
\small
\begin{tabular}{@{}lccccc@{}}
\toprule
\textbf{Prompt} & \textbf{$n$} & \textbf{Latency (ms)} & \textbf{GPU E (J)} & \textbf{ADE (m)} & \textbf{FDE (m)} \\
\midrule
A: straight & 29 & 5897.2 & 193.4 & 0.104 & 0.151 \\
B: right    & 33 & 5074.1 & 185.4 & 0.106 & 0.161 \\
C: stop     & 34 & 4847.0 & 183.9 & 0.082 & 0.139 \\
\midrule
\multicolumn{6}{@{}l}{\textit{Pairwise trajectory / endpoint divergence (m):}} \\
\multicolumn{6}{@{}l}{\quad A vs B: 0.005 / 0.014 \qquad A vs C: 0.016 / 0.033 \qquad B vs C: 0.052 / 0.187} \\
\bottomrule
\end{tabular}
\end{table}

\begin{figure}[htbp]
\centering
\includegraphics[width=\linewidth]{nav_abc_bars.png}
\caption{Stamped nav A/B/C: ego ADE/FDE, model latency, and GPU energy. Accuracy
is similar across prompts; the elevated latency/energy is the CFG-nav dual-pass
cost.}
\label{fig:nav}
\end{figure}

%==============================================================================
\section{Status Audit and Remaining Work}
\label{sec:status}

Table~\ref{tab:audit} is the characterization status, originally audited
2026-08-18 and refreshed 2026-09-02 (\ding{51}~done, \ding{119}~partial,
\ding{55}~not found). The 2026-09-02 refresh closed five items that were open at
the first audit: the MCity camera-normalized FLOP recalculation
(Section~\ref{sec:flops-mcity}), the dataset-matched MCity token-cap sweep
(Section~\ref{sec:tokencap-mcity}), the stamped-nav A/B/C
(Section~\ref{sec:nav}), the same-dataset camera-count ablation
(Section~\ref{sec:cameraablation}), and best-of-K $\mathrm{minADE}_K/\mathrm{minFDE}_K$
(Section~\ref{sec:bestofk}). The remaining open items in Table~\ref{tab:todo}
follow the dataset rule: every non-GT characterization should exist for both
Waymo and MCity, and GT-dependent eval stays Waymo-only until a trusted MCity
future-trajectory target is added.

\begin{table}[h]
\centering
\caption{Characterization status audit (2026-08-18, refreshed 2026-09-02).}
\label{tab:audit}
\small
\begin{tabular}{@{}p{8.6cm}c@{}}
\toprule
\textbf{Characterization item} & \textbf{Status} \\
\midrule
Latency: stage, frame, dataset, token-cap, diffusion & partial \\
Nsight profiling (CUDA API, kernels, stage bottlenecks) & done \\
Operation count / achieved TFLOP/s (camera-normalized) & done \\
Camera-count analysis (controlled same-dataset ablation) & done \\
Energy \& power: avg/peak/per-inference/datasets/sweeps & partial \\
Accuracy: ADE/FDE + token/diffusion/camera effects + best-of-K & done \\
Dataset-matched diffusion-step sweep & done \\
Trajectory-sample (K) replay + best-of-K minADE/minFDE & done \\
Precision BF16/FP16 default-parameter replay & done \\
Precision FP8/INT8 end-to-end replay & not done \\
KV-cache and GPU memory & done \\
Sustained throughput / FPS / drops / input age & done \\
Fixed input configuration documentation & done \\
Token-limit behavior & done \\
Stamped navigation-text A/B/C & done \\
\bottomrule
\end{tabular}
\end{table}

\begin{table}[h]
\centering
\caption{Prioritized remaining work.}
\label{tab:todo}
\small
\begin{tabular}{@{}clcc@{}}
\toprule
\textbf{\#} & \textbf{Item} & \textbf{Waymo} & \textbf{MCity} \\
\midrule
1 & MCity camera-normalized FLOP/TOPS recalculation & done & done \\
2 & Dataset-matched token-cap sweep & done & done \\
3 & Stamped-nav A/B/C at identical settings & done & not done \\
4 & Same-dataset camera-count ablation & done & n/a \\
5 & All-samples $\mathrm{minADE}_K/\mathrm{minFDE}_K$ for K-sweep & done & n/a \\
\midrule
\multicolumn{4}{@{}l}{\textit{Still open:}} \\
6 & MCity accuracy target decision (GT/odometry) & n/a & not done \\
7 & FP8/INT8 end-to-end precision runtime + replay & not done & not done \\
8 & Energy per stage (NVTX-aligned integration) & not done & not done \\
9 & Best-of-K sample selection / ranker in the node & not done & n/a \\
\bottomrule
\end{tabular}
\end{table}

%==============================================================================
\section{Conclusion}
\label{sec:conclusion}

This characterization establishes a reproducible, artifact-backed picture of
Alpamayo~1.5 as a deployed real-time planner on Jetson~AGX~Thor. Four findings
stand out for anyone tuning the system for deployment:
\begin{enumerate}
  \item \textbf{The system is memory-bound in autoregressive decode.} Reason
  decode is $71\%$ of GPU time and the model achieves only $\approx$16\,TFLOP/s,
  so latency tracks decode length, not raw FLOPs.
  \item \textbf{Trajectory-sample count (K) and precision are the expensive
  levers.} K scales latency/energy roughly linearly; FP16 is both slower and
  higher-energy than the BF16 default on this stack, with no accuracy benefit.
  \item \textbf{The reasoning-token cap is nearly inert} on these segments because
  generated length never approaches the cap---diffusion steps and K are the real
  cost knobs.
  \item \textbf{Accuracy is stable and safe at the sanity level.} Ego ADE/FDE
  stays sub-metre across all runnable Waymo configurations and the GT actor check
  finds zero collisions, though these are sanity/clearance checks, not the
  official planning benchmark.
  \item \textbf{Camera count is the strongest accuracy lever---and it is
  expensive.} The controlled same-dataset ablation shows 1$\to$4 cameras cuts ADE
  $3.6\times$ and FDE $6\times$ while costing $2.9\times$ GPU energy and
  $3.7\times$ KV; the camera set, not the reasoning knobs, dominates accuracy.
  \item \textbf{The published trajectory is not the model's best one.} Best-of-K
  recovers $40$--$50\%$ lower ADE/FDE than the published sample~0, so a sample
  selector or ranker is a high-value, no-retrain accuracy win. Conversely,
  stamped nav text barely steers the planner on a scene-constrained segment
  ($\le$\SI{0.19}{\metre} endpoint change) while the CFG dual pass costs
  $\sim$1.6$\times$ latency.
\end{enumerate}
The most valuable open work is now an end-to-end FP8/INT8 runtime (the only
precision path that could plausibly beat BF16), stage-level energy via
NVTX-aligned power integration, a best-of-K sample selector/ranker in the node to
capture the accuracy the model already produces, and a trusted MCity
future-trajectory target so that MCity can move from N/A to real accuracy
reporting.

\end{document}
